\section{Conclusion}\label{sec:conclusion}

Ces trois TPs montrent qu'un m\^{e}me probl\`{e}me algorithmique peut se traduire
par des comportements pratiques tr\`{e}s diff\'{e}rents selon la m\'{e}thode choisie.
Le TP1 met en \'{e}vidence des croissances exponentielle, lin\'{e}aire,
logarithmique et quasi lin\'{e}aire sur des exemples classiques.
Le TP2 illustre
le co\^{u}t tr\`{e}s important d'une exploration exhaustive, m\^{e}me lorsque le
nombre d'entr\'{e}es reste modeste.
Le TP3 confirme enfin l'int\'{e}r\^{e}t des
algorithmes en \(n \log n\) face aux tris quadratiques lorsque la taille des
tableaux augmente.

Sur le plan logiciel, le projet a \'{e}t\'{e} maintenu volontairement simple : un
socle commun minimal pour les mesures et les figures, puis une impl\'{e}mentation
sp\'{e}cifique \`{a} chaque sujet.
Cette approche reste suffisante pour r\'{e}pondre au sujet
sans sur-ing\'{e}nierie, tout en conservant un code typ\'{e}, test\'{e} et facile
\`{a} ex\'{e}cuter.
